% ══════════════════════════════════════════════════════════════════
% Section 8 — Discussion
% ══════════════════════════════════════════════════════════════════
\section{Discussion}
\label{sec:discussion}

\subsection{Architecture trade-offs}
\label{sec:discussion:tradeoffs}

The benchmark results reveal a trade-off between thermodynamic guarantees and
approximation flexibility:

\begin{itemize}
  \item \textbf{MLP} achieves the highest $R^2$ in most cases due to its
    unconstrained weight space.  However, it offers no convexity guarantees,
    which may lead to non-physical behavior outside the training domain and
    convergence issues in implicit solvers.

  \item \textbf{ICNN} guarantees convexity of the energy with respect to the
    invariants, providing material stability by construction.  The
    non-negative weight constraint and skip connections introduce some loss in
    expressiveness, particularly for materials with non-monotonic
    energy--invariant relationships.

  \item \textbf{Polyconvex ICNN} provides the strongest thermodynamic
    guarantee (polyconvexity) at the cost of further restricting the
    approximation space.  The branch-wise decomposition limits cross-invariant
    interactions, which may reduce accuracy for materials where the energy
    depends on invariant products.
\end{itemize}

% TODO: Expand with specific observations from the benchmark data, e.g.:
% - Which materials show the largest gap between MLP and ICNN?
% - Does the polyconvex constraint matter for the materials tested?
% - Are there cases where ICNN outperforms MLP (extrapolation)?

\subsection{Energy--stress loss effectiveness}
\label{sec:discussion:loss}

The dual energy--stress loss (\Cref{eq:energy_stress_loss}) is critical for
obtaining accurate stress predictions from energy-output architectures.
Training on energy alone (setting $\beta = 0$) typically yields
$R^2(\text{stress}) < 0.9$ even when $R^2(\text{energy}) > 0.99$, because
small energy offsets do not penalize gradient errors.

For the ICNN and polyconvex architectures, this loss formulation is the
\emph{only} viable training strategy, since these models output a scalar
energy and cannot be trained on stress targets directly.

% TODO: Consider adding an ablation study (alpha/beta sweep) if space permits.

\subsection{Anisotropic materials}
\label{sec:discussion:anisotropic}

Anisotropic materials (HO, GOH) present additional challenges:
\begin{itemize}
  \item \textbf{Exponential overflow:} The fiber terms involve
    $\exp(b_f\langle I_4 - 1\rangle^2)$, which overflows for large stretches.
    We mitigated this by using reduced parameters ($b_f = 4$ instead of the
    original $b_f = 16.026$) and a tighter deformation range.
  \item \textbf{Higher input dimension:} The 5D invariant input
    ($\bar{I}_1, \bar{I}_2, J, I_4, I_5$) requires more training data and
    larger networks for the same accuracy level.
  \item \textbf{Data generation:} The standard \texttt{create\_datasets}
    pipeline does not support anisotropic materials directly because the
    symbolic SEF involves fiber pseudo-invariants.  We implemented a custom
    pipeline using \texttt{sef\_from\_all\_invariants} with SymPy
    lambdification.
\end{itemize}

% TODO: Discuss whether HO and GOH having the same 5D input but different
% SEFs leads to measurable accuracy differences.

\subsection{Fortran export considerations}
\label{sec:discussion:fortran}

The automated Fortran export eliminates the error-prone manual translation of
NN weights into solver-compatible code.  Key observations:
\begin{itemize}
  \item The embedded-weight approach (compile-time \texttt{PARAMETER} arrays)
    produces a single \texttt{.f90} file with no external dependencies,
    simplifying deployment.
  \item The analytical backpropagation for the Hessian (second derivatives) is
    exact, ensuring quadratic convergence of Newton--Raphson iterations.
  \item The Jaumann correction term in the tangent operator is essential for
    Abaqus compatibility (which uses the Jaumann rate of Cauchy stress).
\end{itemize}

% TODO: Discuss code size, compile time, and runtime overhead of the generated
% Fortran code versus the analytical UMAT.

\subsection{Limitations}
\label{sec:discussion:limitations}

\begin{itemize}
  \item \textbf{CANN excluded.}  The constitutive ANN architecture was
    excluded from the general benchmarks because its exponential basis
    functions ($\exp(b_k \xi^2)$) overflow with standardized inputs.
    Material-specific tuning (see \texttt{examples/model\_discovery\_cann.py})
    is required.
  \item \textbf{Single fiber family.}  The current anisotropic benchmarks
    consider only one fiber direction.  Extension to two-fiber models (e.g.,
    \texttt{HolzapfelOgdenBiaxial}) increases the input dimension to 9 and
    requires further investigation.
  \item \textbf{Training data coverage.}  Synthetic deformation sampling may
    not cover the full deformation space encountered in complex FE
    simulations.  Adaptive sampling strategies could improve robustness.
  \item \textbf{CPU-only timing.}  The inference speed comparison uses
    single-threaded CPU evaluation.  GPU-accelerated inference or batched
    evaluation across integration points is not considered.
  \item \textbf{Single-element validation.}  The FE tests use a single
    element under homogeneous deformation.  Multi-element benchmarks with
    heterogeneous stress fields are left for future work.
\end{itemize}
